% GENERATED FILE. Edit canonical lessons, then run npm run generate:books.
% canonical-source-hash: fnv1a64:8f93b16a389ae784
% canonical-lessons: UR-C511-tarkib, UR-C511-dastarkhvan, UR-C511-pakvana, UR-C511-phunkna, UR-C511-chakhana

\chapter{Laying the Cloth}
\label{ch:ur-laying-the-cloth}
\hlchaptermodality{511}

\begin{chapteropening}
\textbf{By the end of this chapter:} \emph{I can talk about recipes and serving a meal.}

Complete the last lesson of chapter 511: I can talk about recipes and serving a meal.
\end{chapteropening}

\section[\texorpdfstring{\ur{ترکیب} (tarkīb)}{ترکیب (tarkīb)}]{\ur{ترکیب} (tarkīb) — a recipe; a method}

\label{lesson:UR-C511-tarkib}



\begin{quote}
Before the new one: say the Urdu for to have (something) ground, then the Urdu for to squeeze out.
\end{quote}

\subsection*{You'll want to know: \ur{ترکیب}}
\textbf{\ur{ترکیب}} — \emph{tarkīb} — \textquotedblleft{}a recipe; a method\textquotedblright{}.

\begin{quote}
\ur{یہ} \ur{نانی} \ur{کی} \ur{ترکیب} \ur{ہے۔}
\end{quote}

\begin{quote}
\emph{yih nānī kī tarkīb hai.}
\end{quote}

\begin{quote}
This is grandmother's recipe.
\end{quote}

Feminine.

\begin{cousinweb}
\textbf{\ur{ترکیب}} (\emph{tarkīb}) is usually traced, through Persian, to Arabic \emph{tarkīb}, \textquotedblleft{}putting together\textquotedblright{}.
\end{cousinweb}

\begin{grammarlens}[title={What you've built}]
One more word for describing what you lived and felt.
\end{grammarlens}

\subsection*{Your turn}
\begin{itemize}
\raggedright
  \item \emph{Say it:} \emph{tarkīb}
  \item \emph{Say it:} \emph{tarkīb}, once more
  \item \emph{Say it:} \emph{yih nānī kī tarkīb hai.} — This is grandmother's recipe.
  \item \emph{Recall:} say the Urdu for to have (something) ground, then the Urdu for to squeeze out, then say \emph{tarkīb} again
\end{itemize}

\begin{tcolorbox}[breakable,colback=teal!4,colframe=teal!35!black,title={Before you move on}]
What does \ur{ترکیب} mean? (\textquotedblleft{}A recipe; a method\textquotedblright{}.) Say it once more.
\end{tcolorbox}
\section[\texorpdfstring{\ur{دسترخوان} (dastarkhvān)}{دسترخوان (dastarkhvān)}]{\ur{دسترخوان} (dastarkhvān) — a cloth spread for a meal}

\label{lesson:UR-C511-dastarkhvan}



\begin{quote}
Before the new one: say the Urdu for to squeeze out, then the Urdu for a recipe.
\end{quote}

\subsection*{You'll want to know: \ur{دسترخوان}}
\textbf{\ur{دسترخوان}} — \emph{dastarkhvān} — \textquotedblleft{}a cloth spread for a meal\textquotedblright{}.

\begin{quote}
\ur{امی} \ur{نے} \ur{فرش} \ur{پر} \ur{دسترخوان} \ur{بچھایا۔}
\end{quote}

\begin{quote}
\emph{ammī ne farsh par dastarkhvān bichhāyā.}
\end{quote}

\begin{quote}
Mum spread the dining cloth on the floor.
\end{quote}

Masculine.

\begin{cousinweb}
\textbf{\ur{دسترخوان}} (\emph{dastarkhvān}) is usually traced to Persian \emph{dastarkhwān}, \textquotedblleft{}a cloth spread for a meal\textquotedblright{}.
\end{cousinweb}

\begin{grammarlens}[title={What you've built}]
One more word for describing what you lived and felt.
\end{grammarlens}

\subsection*{Your turn}
\begin{itemize}
\raggedright
  \item \emph{Say it:} \emph{dastarkhvān}
  \item \emph{Say it:} \emph{dastarkhvān}, once more
  \item \emph{Say it:} \emph{ammī ne farsh par dastarkhvān bichhāyā.} — Mum spread the dining cloth on the floor.
  \item \emph{Recall:} say the Urdu for to squeeze out, then the Urdu for a recipe, then say \emph{dastarkhvān} again
\end{itemize}

\begin{tcolorbox}[breakable,colback=teal!4,colframe=teal!35!black,title={Before you move on}]
What does \ur{دسترخوان} mean? (\textquotedblleft{}A cloth spread for a meal\textquotedblright{}.) Say it once more.
\end{tcolorbox}
\section[\texorpdfstring{\ur{پکوانا} (pakvānā)}{پکوانا (pakvānā)}]{\ur{پکوانا} (pakvānā) — to have (food) cooked}

\label{lesson:UR-C511-pakvana}



\begin{quote}
Before the new one: say the Urdu for a recipe, then the Urdu for a cloth spread for a meal.
\end{quote}

\subsection*{You'll want to know: \ur{پکوانا}}
\textbf{\ur{پکوانا}} — \emph{pakvānā} — \textquotedblleft{}to have (food) cooked\textquotedblright{}.

\begin{quote}
\ur{ابو} \ur{نے} \ur{مہمانوں} \ur{کے} \ur{لیے} \ur{بریانی} \ur{پکوائی۔}
\end{quote}

\begin{quote}
\emph{abbū ne mehmānoṅ ke liye biryānī pakvāī.}
\end{quote}

\begin{quote}
Dad had biryani cooked for the guests.
\end{quote}

Past, with \ur{نے} (he or she …): \ur{اس} \ur{نے} \ur{حلوہ} \ur{پکوایا}, \emph{us ne halvā pakvāyā}.

\begin{cousinweb}
\textbf{\ur{پکوانا}} (\emph{pakvānā}) is the second causative of \emph{paknā}, \textquotedblleft{}to be cooked\textquotedblright{}, which CDIAL lists under Sanskrit \emph{pakva}, \textquotedblleft{}cooked, ripe\textquotedblright{} (CDIAL 7621): to have someone cook.
\end{cousinweb}

\begin{grammarlens}[title={What you've built}]
One more word for describing what you lived and felt.
\end{grammarlens}

\subsection*{Your turn}
\begin{itemize}
\raggedright
  \item \emph{Say it:} \emph{pakvānā}
  \item \emph{Say it:} \emph{pakvānā}, once more
  \item \emph{Say it:} \emph{abbū ne mehmānoṅ ke liye biryānī pakvāī.} — Dad had biryani cooked for the guests.
  \item \emph{Recall:} say the Urdu for a recipe, then the Urdu for a cloth spread for a meal, then say \emph{pakvānā} again
\end{itemize}

\begin{tcolorbox}[breakable,colback=teal!4,colframe=teal!35!black,title={Before you move on}]
What does \ur{پکوانا} mean? (\textquotedblleft{}To have (food) cooked\textquotedblright{}.) Say it once more.
\end{tcolorbox}
\section[\texorpdfstring{\ur{پھونکنا} (phūṅknā)}{پھونکنا (phūṅknā)}]{\ur{پھونکنا} (phūṅknā) — to blow (on)}

\label{lesson:UR-C511-phunkna}



\begin{quote}
Before the new one: say the Urdu for a cloth spread for a meal, then the Urdu for to have (food) cooked.
\end{quote}

\subsection*{You'll want to know: \ur{پھونکنا}}
\textbf{\ur{پھونکنا}} — \emph{phūṅknā} — \textquotedblleft{}to blow (on)\textquotedblright{}.

\begin{quote}
\ur{بچہ} \ur{گرم} \ur{دودھ} \ur{پھونک} \ur{کر} \ur{پیتا} \ur{ہے۔}
\end{quote}

\begin{quote}
\emph{bachchā garm dūdh phūṅk kar pītā hai.}
\end{quote}

\begin{quote}
The child blows on hot milk before drinking it.
\end{quote}

Past, with \ur{نے} (he or she …): \ur{اس} \ur{نے} \ur{چولہے} \ur{کی} \ur{آگ} \ur{پھونکی}, \emph{us ne chūlhe kī āg phūṅkī}.

\begin{cousinweb}
\textbf{\ur{پھونکنا}} (\emph{phūṅknā}) is inherited: CDIAL lists the Hindi cousin \emph{phū̃knā}, \textquotedblleft{}to blow\textquotedblright{}, under \emph{*phūtka}, a form it reconstructs (CDIAL 9102).
\end{cousinweb}

\begin{grammarlens}[title={What you've built}]
One more word for describing what you lived and felt.
\end{grammarlens}

\subsection*{Your turn}
\begin{itemize}
\raggedright
  \item \emph{Say it:} \emph{phūṅknā}
  \item \emph{Say it:} \emph{phūṅknā}, once more
  \item \emph{Say it:} \emph{bachchā garm dūdh phūṅk kar pītā hai.} — The child blows on hot milk before drinking it.
  \item \emph{Recall:} say the Urdu for a cloth spread for a meal, then the Urdu for to have (food) cooked, then say \emph{phūṅknā} again
\end{itemize}

\begin{tcolorbox}[breakable,colback=teal!4,colframe=teal!35!black,title={Before you move on}]
What does \ur{پھونکنا} mean? (\textquotedblleft{}To blow (on)\textquotedblright{}.) Say it once more.
\end{tcolorbox}
\section[\texorpdfstring{\ur{چکھانا} (chakhānā)}{چکھانا (chakhānā)}]{\ur{چکھانا} (chakhānā) — to give (someone) a taste}

\label{lesson:UR-C511-chakhana}



\begin{quote}
Before the new one: say the Urdu for to have (food) cooked, then the Urdu for to blow (on).
\end{quote}

\subsection*{You'll want to know: \ur{چکھانا}}
\textbf{\ur{چکھانا}} — \emph{chakhānā} — \textquotedblleft{}to give (someone) a taste\textquotedblright{}.

\begin{quote}
\ur{امی} \ur{نے} \ur{مہمان} \ur{کو} \ur{حلوہ} \ur{چکھایا۔}
\end{quote}

\begin{quote}
\emph{ammī ne mehmān ko halvā chakhāyā.}
\end{quote}

\begin{quote}
Mum gave the guest a taste of the halwa.
\end{quote}

Past, with \ur{نے} (he or she …): \ur{اس} \ur{نے} \ur{مہمانوں} \ur{کو} \ur{حلوہ} \ur{چکھایا}, \emph{us ne mehmānoṅ ko halvā chakhāyā}.

\begin{cousinweb}
\textbf{\ur{چکھانا}} (\emph{chakhānā}) is the causative of \emph{chakhnā}, \textquotedblleft{}to taste\textquotedblright{}; CDIAL lists the Hindi cousin \emph{cākhnā}, \textquotedblleft{}to taste\textquotedblright{}, under \emph{*cakṣati}, \textquotedblleft{}tastes\textquotedblright{}, a form it reconstructs (CDIAL 4557).
\end{cousinweb}

\begin{grammarlens}[title={What you've built}]
One more word for describing what you lived and felt.
\end{grammarlens}

\subsection*{Your turn}
\begin{itemize}
\raggedright
  \item \emph{Say it:} \emph{chakhānā}
  \item \emph{Say it:} \emph{chakhānā}, once more
  \item \emph{Say it:} \emph{ammī ne mehmān ko halvā chakhāyā.} — Mum gave the guest a taste of the halwa.
  \item \emph{Recall:} say the Urdu for to have (food) cooked, then the Urdu for to blow (on), then say \emph{chakhānā} again
\end{itemize}

\begin{tcolorbox}[breakable,colback=teal!4,colframe=teal!35!black,title={Before you move on}]
What does \ur{چکھانا} mean? (\textquotedblleft{}To give (someone) a taste\textquotedblright{}.) Say it once more.
\end{tcolorbox}
